\chapter{Detecting entanglement: the Peres criterion}\label{ch:peres}

\section{Separable mixed states}
For pure states, entanglement means ``not a product'' (\cref{prop:det}). For mixed states the definition must allow for classical correlations.
\begin{definition}[Werner, 1989]
A density matrix of $AB$ is \emph{separable} if it is a convex combination of product states,
\begin{equation}\label{eq:separable}
  \hat\rho=\sum_ip_i\,\hat\rho_i^A\otimes\hat\rho_i^B,\qquad p_i\ge0,\ \sum_ip_i=1.
\end{equation}
Otherwise it is \emph{entangled}.
\end{definition}
Separable states are exactly those that Alice and Bob can prepare with local operations and classical communication (LOCC): they agree on $i$ with probability $p_i$
(e.g.\ by a phone call) and prepare $\hat\rho_i^A$ and $\hat\rho_i^B$ locally. For example $\frac12(\proj{00}+\proj{11})$ is correlated (both qubits are always equal in the
$Z$ basis) but separable. In contrast, $\proj{\Phi^+}$ has the same diagonal but also the coherences $\ketbra{00}{11}$, which make it entangled: it is correlated in
every basis. Deciding whether a given $\hat\rho$ is separable is hard in general (NP-hard in the dimension, Gurvits 2003), but there are simple necessary criteria.

\section{Partial transpose}
Let $\hat\rho$ be a density matrix on $\Ha_A\otimes\Ha_B$:
\begin{equation}
  \hat\rho=\sum_{ijkl}p^{ij}_{kl}\ketbra ij\otimes\ketbra kl.
\end{equation}
Its \emph{partial transpose} with respect to $B$ is
\begin{equation}
  \hat\rho^{T_B}=(I\otimes T)\hat\rho=\sum_{ijkl}p^{ij}_{kl}\ketbra ij\otimes\left(\ketbra kl\right)^T=\sum_{ijkl}p^{ij}_{kl}\ketbra ij\otimes\ketbra lk
  =\sum_{ijkl}p^{ij}_{lk}\ketbra ij\otimes\ketbra kl,
\end{equation}
where in the last step the summation indices $k$ and $l$ have been exchanged. In terms of matrix elements, only the indices of subsystem $B$ are exchanged:
\begin{equation}
  \left(\rho^{T_B}\right)_{m\mu,n\nu}=\rho_{m\nu,n\mu},
\end{equation}
where Latin indices refer to $A$ and Greek indices to $B$. With the big-endian ordering, $\hat\rho$ is a block matrix made of $\dim\Ha_A\times\dim\Ha_A$ blocks, each of size
$\dim\Ha_B\times\dim\Ha_B$:
\begin{equation}
  \hat\rho=\begin{pmatrix}A_{11}&A_{12}&\cdots\\A_{21}&A_{22}&\cdots\\ \vdots&\vdots&\ddots\end{pmatrix}\quad\Longrightarrow\quad
  \hat\rho^{T_B}=\begin{pmatrix}A_{11}^T&A_{12}^T&\cdots\\A_{21}^T&A_{22}^T&\cdots\\ \vdots&\vdots&\ddots\end{pmatrix}.
\end{equation}
Each block is transposed, but the order of the blocks is not changed. (The partial transpose depends on the basis, but its eigenvalues do not.)

\section{The Peres--Horodecki criterion}
\begin{theorem}[Peres \cite{Peres1996}; Horodecki \cite{Horodecki1996}]
If $\hat\rho$ is separable, then all the eigenvalues of $\hat\rho^{T_B}$ are non-negative (``PPT'': positive partial transpose). Equivalently, if $\hat\rho^{T_B}$ has a negative
eigenvalue, then $\hat\rho$ is entangled. For systems of dimension $2\times2$ and $2\times3$, the converse is also true: PPT implies separability.
\end{theorem}
\begin{proof}[Proof of the first part]
A separable state is a convex sum of product states, \cref{eq:separable}. Then
\begin{equation}
  \hat\rho^{T_B}=(I\otimes T)\hat\rho=\sum_ip_i\,\hat\rho_i^A\otimes\left(\hat\rho_i^B\right)^T.
\end{equation}
Transposition preserves the eigenvalues, so $(\hat\rho_i^B)^T$ has the same non-negative eigenvalues as $\hat\rho_i^B$, and is itself a valid density matrix (it is
also the complex conjugate of $\hat\rho_i^B$, since $\hat\rho_i^B$ is Hermitian). $\hat\rho^{T_B}$ is therefore a convex sum of positive semi-definite operators, and is positive
semi-definite. A negative eigenvalue violates this requirement.
\end{proof}
\begin{remark}[Why the transpose?]
The transpose $T$ is a \emph{positive} map (it maps density matrices to density matrices) which is not \emph{completely positive}: $I\otimes T$ applied to an entangled
state can give a non-positive operator. Physically, $\hat\rho^T=\hat\rho^*$ is the time-reversed state; reversing time on one half of an entangled system only is not a
physical operation. Every positive but not completely positive map gives an entanglement criterion, and the set of all of them characterizes separability (Horodecki 1996).
In higher dimensions there exist entangled PPT states (``bound entanglement''), which cannot be distilled into Bell pairs.
\end{remark}

\section{Examples}
\begin{example}[An entangled state]
For $\ket\psi=(\ket{01}+\ket{10})/\sqrt2$, in the basis $\{\ket{00},\ket{01},\ket{10},\ket{11}\}$,
\begin{equation}
  \hat\rho=\frac12\begin{pmatrix}0&0&0&0\\0&1&1&0\\0&1&1&0\\0&0&0&0\end{pmatrix},
\end{equation}
where, for example, $\rho_{10,01}=\braket{10}\psi\braket\psi{01}=\frac12$. The partial transpose moves the off-diagonal elements: $\rho_{01,10}\to(\rho^{T_B})_{00,11}$ and
$\rho_{10,01}\to(\rho^{T_B})_{11,00}$ (the upper-right block $\begin{pmatrix}0&0\\1&0\end{pmatrix}$ becomes $\begin{pmatrix}0&1\\0&0\end{pmatrix}$), so that
\begin{equation}
  \hat\rho^{T_B}=\frac12\begin{pmatrix}0&0&0&1\\0&1&0&0\\0&0&1&0\\1&0&0&0\end{pmatrix}.
\end{equation}
Its eigenvalues follow from (omitting the factor $\frac12$)
\begin{equation}
  \det\begin{pmatrix}-\lambda&0&0&1\\0&1-\lambda&0&0\\0&0&1-\lambda&0\\1&0&0&-\lambda\end{pmatrix}=\lambda^2(1-\lambda)^2-(1-\lambda)^2=(1-\lambda)^2(\lambda-1)(\lambda+1)=0.
\end{equation}
The eigenvalues of $\hat\rho^{T_B}$ are therefore $\frac12,\frac12,\frac12$ and $\lambda_4=-\frac12<0$: the state is entangled. The eigenvector of the negative eigenvalue is
$(\ket{00}-\ket{11})/\sqrt2$.
\end{example}

\begin{example}[A separable state]
For $\ket{\psi_2}=(\ket{10}+\ket{00})/\sqrt2=\frac1{\sqrt2}(\ket0+\ket1)\otimes\ket0$,
\begin{equation}
  \hat\rho=\frac12\begin{pmatrix}1&0&1&0\\0&0&0&0\\1&0&1&0\\0&0&0&0\end{pmatrix}\quad\Longrightarrow\quad\hat\rho^{T_B}=\hat\rho,
\end{equation}
because every block is symmetric. The partial transpose is the same matrix, with eigenvalues $1,0,0,0\ge0$, as expected for a separable state.
\end{example}

\begin{example}[The Werner state]\label{ex:werner}
Consider the mixture of a singlet with white noise, $\hat\rho_W=p\proj{\psi^-}+(1-p)\frac I4$, $0\le p\le1$. This models a Bell pair sent through a depolarizing channel.
Since $I^{T_B}=I$ and $\proj{\psi^-}^{T_B}$ has eigenvalues $\frac12,\frac12,\frac12,-\frac12$ (as in the first example, up to a local unitary), $\hat\rho_W^{T_B}$ has the eigenvalues
\begin{equation}
  \frac p2+\frac{1-p}4\ (\times3),\qquad-\frac p2+\frac{1-p}4=\frac{1-3p}4.
\end{equation}
The Werner state is entangled if and only if $p>1/3$. (It violates the CHSH inequality only for $p>1/\sqrt2$: entanglement and nonlocality are not the same thing.)
\end{example}

\section{Quantifying entanglement}
\begin{itemize}
  \item The \emph{negativity} $\mathcal N=\sum_{\lambda_i<0}|\lambda_i(\hat\rho^{T_B})|$ ($=\frac12$ for a Bell state, $\frac{3p-1}4$ for the Werner state with $p>1/3$) and the
  logarithmic negativity $E_N=\log_2(1+2\mathcal N)$ are easily computable entanglement measures.
  \item The \emph{concurrence} of Wootters (1998), $C(\hat\rho)=\max(0,\lambda_1-\lambda_2-\lambda_3-\lambda_4)$, where $\lambda_i$ are the square roots of the eigenvalues of
  $\hat\rho(Y\otimes Y)\hat\rho^*(Y\otimes Y)$ in decreasing order, gives the exact entanglement of formation of any two-qubit state. For pure states it reduces to
  $2|\alpha\delta-\beta\gamma|$ (\cref{prop:det}).
  \item An \emph{entanglement witness} is an observable $W$ with $\tr(W\hat\rho)\ge0$ for all separable states. Measuring $\tr(W\hat\rho)<0$ proves entanglement with a single
  observable, without full tomography. For a target state $\ket\psi$, $W=\alpha I-\proj\psi$ with $\alpha$ the maximal squared overlap of $\ket\psi$ with product states ($\frac12$
  for Bell states) is a witness: a fidelity with a Bell state larger than $1/2$ proves entanglement.
\end{itemize}

\section{Exercises}
\begin{exercise}
Show that the partial transpose with respect to $A$ has the same spectrum as the partial transpose with respect to $B$. (Hint: $\hat\rho^{T_A}=(\hat\rho^{T_B})^T$.)
\end{exercise}
\begin{exercise}
For which values of $p$ is $\hat\rho=p\proj{\Phi^+}+(1-p)\proj{01}$ entangled?
(Answer: for all $p>0$. Compute the eigenvalue of the block $\{\ket{01},\ket{10}\}$ of $\hat\rho^{T_B}$.)
\end{exercise}
